\begin{center}
    \textbf{\Huge{Visão de Produto}}
\end{center}

\section{Introdução}

O propósito deste documento é coletar, analisar e definir as necessidades de alto nível e as principais \emph{features} do ASA$^{2}$/MIIA. O documento foca nas capacidades necessárias aos \emph{stakeholders} e usuários-alvo, bem como nas motivações que justificam essas necessidades. Os detalhes de como o ASA$^{2}$/MIIA atende tais necessidades são especificados em artefatos complementares, incluindo casos de uso, \emph{user stories} e documentação suplementar.

\subsection{Referências}

\begin{itemize}
    \item Documento Institucional Preliminar, ``\textbf{Ambiente de Simulação Aeroespacial Avançado -- ASA$^{2}$}'', Instituto de Estudos Avançados (IEAv), 2025.
\end{itemize}

\newpage

\section{Posicionamento}

\subsection{Declaração do Problema}

\begin{table}[h]
    \begin{tabular}{|p{0.3\textwidth}|p{0.65\textwidth}|}
        \hline
            \textbf{O problema de} &
            Integrar e executar, de forma determinística e com baixa latência, agentes de decisão baseados em modelos de inteligência artificial associados a simuladores construtivos e virtuais, garantindo estabilidade, previsibilidade temporal e isolamento operacional. \\
        \hline
            \textbf{afeta} &
            Engenheiros de simulação, integradores de sistemas, equipes de P\&D e operadores que dependem de agentes autônomos ou semi-autônomos capazes de atuar em ciclos de simulação de alta frequência. \\
        \hline
            \textbf{cujo impacto é} &
            A impossibilidade de empregar agentes de IA em simulações de nível operacional; integrações frágeis, não determinísticas ou de difícil depuração; instabilidade do simulador; aumento de retrabalho e dívida técnica; e a incapacidade de validar cenários com rigor metodológico. \\
        \hline
            \textbf{uma solução bem-sucedida seria} &
            Um mecanismo padronizado, documentado e tecnicamente validado que permita executar agentes de decisão com desempenho consistente, latência controlada, comportamento reproduzível e um processo de integração simples, rastreável e confiável. \\
        \hline
    \end{tabular}
\end{table}

A solução envolve um conjunto de propriedades complementares:

\begin{itemize}

    \item \textbf{Determinismo:} para um mesmo conjunto de observações e estado interno, o agente deve produzir exatamente a mesma ação, de forma consistente entre execuções. Isso elimina variabilidade indesejada e permite a reprodução rigorosa de cenários.

    \item \textbf{Previsibilidade temporal:} o tempo necessário para que o agente processe suas entradas e gere uma ação deve permanecer dentro de limites bem definidos. A simulação depende de ciclos estáveis; qualquer variação temporal pode comprometer a integridade do ambiente virtual.

    \item \textbf{Baixa latência:} a inferência deve ser suficientemente rápida para não atrasar o avanço do simulador. Em simulações de alta frequência, atrasos acumulados inviabilizam a operação.

    \item \textbf{Isolamento operacional:} falhas, exceções ou comportamentos inesperados do agente não podem comprometer o simulador. A solução deve permitir detectar, conter e registrar falhas sem interromper a execução global.

    \item \textbf{Reprodutibilidade:} o comportamento observado de um agente deve poder ser repetido em diferentes execuções, máquinas ou versões do ambiente, desde que os mesmos artefatos (modelo, parâmetros e \emph{seeds}) sejam utilizados.

    \item \textbf{Rastreabilidade e verificação:} cada agente deve possuir metadados claros (versão, origem, parâmetros e dependências). A execução deve permitir inspeção, instrumentação e validação contínua, facilitando a depuração e a avaliação metodológica.

\end{itemize}

Em conjunto, essas propriedades permitem a utilização confiável de agentes de IA em cenários operacionais, experimentais e de validação, preservando a integridade temporal e operacional do ciclo de simulação.

\subsection{Declaração de Posicionamento do Produto}

\begin{table}[h]
    \begin{tabular}{|p{0.2\textwidth}|p{0.75\textwidth}|}
        \hline
            \textbf{Para} &
            Engenheiros de simulação, integradores de sistemas e equipes de P\&D que precisam empregar agentes de inteligência artificial em simuladores construtivos e virtuais. \\
        \hline
            \textbf{Que} &
            Enfrentam limitações técnicas ao integrar agentes treinados em ambientes externos ao simulador, como latência imprevisível, falta de determinismo, acoplamento excessivo e dificuldades de validação e manutenção. \\
        \hline
            \textbf{O ASA$^{2}$/MIIA} &
            É um mecanismo especializado para integração e execução de agentes de decisão baseados em IA. \\
        \hline
            \textbf{Que} &
            Permite que esses agentes sejam implantados e executados de forma eficiente, determinística, com baixa latência e reprodutibilidade garantida dentro do ciclo de simulação. \\
        \hline
            \textbf{Diferentemente de} &
            Servidores genéricos de inferência, arquiteturas distribuídas ou integrações \emph{ad-hoc} frágeis e não reprodutíveis. \\
        \hline
            \textbf{Nosso produto} &
            Foca especificamente no domínio de simulação construtiva e virtual, priorizando desempenho temporal, estabilidade operacional, previsibilidade do agente, rastreabilidade do modelo e baixo acoplamento com o simulador. \\
        \hline
    \end{tabular}
\end{table}

\subsection{Limites e Não-Objetivos do Produto}

O ASA\textsuperscript{2}/MIIA tem como objetivo fornecer mecanismos controlados e previsíveis para a integração e execução de agentes de IA em simuladores construtivos e virtuais. No entanto, este produto não se propõe a:

\begin{itemize}
    \item Garantir determinismo absoluto independentemente de variações de hardware, compiladores, bibliotecas numéricas ou configurações de execução fora do ambiente controlado.
    \item Substituir simuladores, motores físicos ou \emph{pipelines} de treinamento de agentes de inteligência artificial.
    \item Corrigir deficiências intrínsecas dos modelos de IA, tais como instabilidade de políticas, vieses de treinamento ou comportamentos emergentes indesejados.
    \item Eliminar a necessidade de validação metodológica, testes experimentais e análise crítica dos resultados produzidos pelos agentes.
\end{itemize}

Esses limites são assumidos explicitamente como parte do escopo do produto.


\section{Descrição dos \emph{Stakeholders} e Usuários}

Esta seção descreve os principais \emph{stakeholders} e usuários envolvidos no projeto, bem como os problemas fundamentais percebidos no contexto atual de integração de agentes de inteligência artificial em simulações construtivas e virtuais. O objetivo é fornecer o pano de fundo e a justificativa para a definição dos requisitos do sistema.

\subsection{Resumo dos \emph{Stakeholders}}

\begin{table}[h]
    \begin{tabular}{| p{0.2\textwidth} | p{0.3\textwidth} | p{0.42\textwidth} |}
        \hline Nome & Descrição & Responsabilidades \\ \hline
        
        \textbf{Product Owner (PO)} &
        Responsável pela visão, prioridades e valor do produto &
        Define prioridades, valida entregas e aprova requisitos e escopo \\ \hline
        
        \textbf{Pesquisador de Agentes de IA} &
        Responsável por gerar, treinar e avaliar agentes de decisão &
        Fornece artefatos treinados, valida comportamento e apoia decisões técnicas relacionadas às políticas \\ \hline
        
        \textbf{Engenheiro de Simulação} &
        Responsável pelo simulador construtivo/virtual e por sua manutenção &
        Garante compatibilidade, integra o mecanismo ao simulador e assegura estabilidade temporal \\ \hline

    \end{tabular}
\end{table}

\newpage

\subsection{Resumo dos Usuários}

\begin{table}[h]
    \begin{tabular}{| p{0.2\textwidth} | p{0.2\textwidth} | p{0.3\textwidth} | p{0.2\textwidth} | }
        \hline Nome & Descrição & Responsabilidades & \emph{Stakeholder} \\ \hline
              
        \textbf{Pesquisador de Agentes de IA} &
        Treinador e analista das políticas &
        Gera agentes, interpreta falhas e avalia comportamento e desempenho &
        Pesquisador de Agentes de IA \\ \hline
        
        \textbf{Analista de Simulação} &
        Usuário final do ambiente de simulação &
        Executa cenários, observa os efeitos dos agentes e reporta problemas &
        Product Owner \\ \hline

    \end{tabular}
\end{table}

\subsection{Ambiente de Usuário}

Os usuários operam predominantemente em estações de trabalho Linux, integradas ao ambiente institucional do ASA\textsuperscript{2}.

\subsubsection{Pesquisador de Agentes de IA}

Atua no desenvolvimento e treinamento de agentes utilizando ferramentas baseadas em Python. O processo envolve a criação de ambientes de simulação para treinamento, a instrumentação de experimentos, o ajuste de hiperparâmetros e a geração de artefatos de modelo. Os agentes resultantes são exportados em formatos derivados do ecossistema PyTorch, incluindo metadados adicionais de versão e configuração.

\subsubsection{Analista de Simulação}

Utiliza interfaces de análise baseadas em Python (como JupyterHub ou Jupyter Notebooks) e APIs fornecidas pelo simulador para inspeção de estados, coleta de métricas e geração de relatórios sobre cenários experimentais.

\newpage

\subsection{Resumo das Principais Necessidades dos \emph{Stakeholders} e Usuários}

As necessidades a seguir refletem dois eixos complementares: (i) propriedades críticas de execução dos agentes no ciclo de simulação e (ii) requisitos de integração, governança técnica e manutenção do produto ao longo do tempo.

\begin{table}[h]
    \begin{tabular}{| p{0.20\textwidth} | p{0.12\textwidth} | p{0.15\textwidth} | p{0.18\textwidth} | p{0.21\textwidth} |}
        \hline
        \textbf{Necessidade} & \textbf{Prioridade} & \textbf{Preocupações} & \textbf{Solução Atual} & \textbf{Soluções Propostas} \\ \hline
        
        Execução determinística dos agentes &
        Alta &
        Variabilidade temporal, efeitos estocásticos &
        Não há &
        Mecanismo padronizado com controle explícito de estado e temporização \\ \hline
        
        Previsibilidade temporal no ciclo de simulação &
        Alta &
        Overhead de comunicação, dependência de Python &
        Não há &
        Execução integrada ao ciclo do simulador, adaptada a alta frequência \\ \hline
        
        Reprodutibilidade entre execuções e ambientes &
        Alta &
        Sensibilidade a \emph{seeds} &
        Não há &
        \emph{Pipeline} formal de exportação, validação e inspeção de agentes \\ \hline
        
        Rastreabilidade e versionamento dos agentes &
        Média &
        Inconsistência de artefatos, perda de histórico &
        Não há &
        Processo padronizado de metadados, versionamento e documentação técnica \\ \hline
        
        Integração simples e de baixo acoplamento &
        Média &
        Dependências cruzadas, impacto no simulador &
        Não há &
        API estável, modular e independente da lógica interna do simulador \\ \hline
        
    \end{tabular}
\end{table}

\subsection{Alternativas e Competição}

No contexto do ASA\textsuperscript{2}, diferentes abordagens técnicas são consideradas para a integração e execução de agentes de inteligência artificial em simuladores construtivos e virtuais. Essas abordagens representam alternativas tecnológicas viáveis, mas com diferentes compromissos em termos de desempenho temporal, determinismo, acoplamento e reprodutibilidade.

As principais classes de alternativas incluem:

\begin{itemize}
    \item \textbf{Runtimes interoperáveis de inferência} (ex.: ONNX Runtime): oferecem portabilidade entre modelos e ferramentas de treinamento, além de bom desempenho em inferência. Entretanto, carecem de mecanismos específicos para controle de estado, sincronização temporal e rastreabilidade no contexto de simulação construtiva, exigindo adaptações adicionais para garantir determinismo e previsibilidade operacional.

    \item \textbf{Runtimes nativos de inferência} (ex.: Libtorch, TensorFlow C++ API): permitem integração direta em aplicações C++ com menor overhead e maior controle do ciclo de execução. Apesar do desempenho favorável, essas abordagens tendem a introduzir forte acoplamento com frameworks específicos, além de maior complexidade de manutenção e validação ao longo do tempo.

    \item \textbf{Execução remota de agentes} (ex.: REST ou gRPC): proporciona flexibilidade arquitetural e isolamento de processos, facilitando a reutilização de agentes treinados externamente. Contudo, a comunicação remota introduz latência adicional, variabilidade temporal e dependência de infraestrutura de rede, tornando essa abordagem inadequada para ciclos de simulação de alta frequência que exigem previsibilidade rigorosa.

    \item \textbf{\emph{Embedding} de Python em aplicações C++}: possibilita acesso direto a modelos e ecossistemas de treinamento amplamente utilizados. No entanto, essa estratégia apresenta elevado acoplamento, dificuldades de depuração, e riscos à estabilidade e à manutenção do simulador.
\end{itemize}

Embora cada uma dessas alternativas apresente vantagens em contextos específicos, nenhuma atende simultaneamente aos requisitos de determinismo, previsibilidade temporal, baixa latência, reprodutibilidade e baixo acoplamento exigidos pelo ambiente de simulação construtiva e virtual do ASA\textsuperscript{2}. O ASA\textsuperscript{2}/MIIA posiciona-se como uma solução especializada que busca endereçar essas limitações de forma integrada e alinhada ao domínio de simulação.

\section{\emph{Overview} do Produto}

\subsection{Perspectiva do Produto}

O ASA\textsuperscript{2}/MIIA é um componente de infraestrutura especializado na integração e execução de agentes de decisão em simuladores construtivos e virtuais. Ele opera diretamente no ciclo de simulação, provendo mecanismos de execução determinística, previsibilidade temporal, rastreabilidade e uma interface padronizada para agentes treinados externamente.

Ele não substitui o simulador nem os \emph{pipelines} de treinamento de IA; atua como uma camada intermediária especializada, responsável por garantir propriedades críticas de execução, alinhada às diretrizes institucionais do ASA\textsuperscript{2} para integração padronizada de modelos de inteligência artificial.

\subsection{Suposições e Dependências}

\begin{itemize}
    \item Disponibilidade de agentes treinados em PyTorch, SB3 ou ferramentas equivalentes.
    \item Execução predominante em ambiente Linux de alto desempenho.
    \item Requisitos operacionais de baixa latência e alta previsibilidade temporal no ciclo do simulador.
    \item Disponibilidade de especialistas técnicos em C++, simulação e aprendizado por reforço.
\end{itemize}

\section{\emph{Features} do Produto}

\begin{itemize}
    \item Execução determinística de agentes de decisão integrados ao simulador.
    \item Suporte a múltiplos formatos de exportação de políticas (ex.: PyTorch, ONNX).
    \item Mecanismos de rastreabilidade, versionamento e metadados dos agentes.
    \item Instrumentação e monitoramento do desempenho do agente durante a simulação.
    \item \emph{Pipeline} de validação mínima do agente (integridade, compatibilidade e \emph{seeds}).
\end{itemize}

\section{Outros Requisitos do Produto}

\begin{itemize}
    \item Latência compatível com ciclos de simulação de alta frequência.
    \item Robustez contra falhas no carregamento, inicialização ou execução de agentes.
    \item API C++ simples, estável e de baixo acoplamento com o simulador.
    \item Documentação técnica cobrindo integração, operação, validação e \emph{troubleshooting}.
    \item Garantias de segurança como requisito estrutural do produto, incluindo mecanismos para prevenir, detectar e conter a execução de agentes maliciosos ou não conformes no contexto da simulação.
\end{itemize}
